\begin{ex}
Biểu đồ cột biểu diễn số cây trồng của 4 khối: Khối 1: 50 cây, Khối 2: 60 cây, Khối 3: 70 cây, Khối 4: 80 cây. Khối nào trồng được nhiều cây nhất? (Khối 4 ghi 4).
\shortans{4}
\loigiai{
Khối 4 trồng 80 cây.
}
\end{ex}

\begin{ex}
Khối 4 trồng nhiều hơn Khối 1 bao nhiêu cây?
\shortans{30}
\loigiai{
$80 - 50 = 30$ cây.
}
\end{ex}

\begin{ex}
Tổng số cây cả 4 khối trồng được là bao nhiêu cây?
\shortans{260}
\loigiai{
$50 + 60 + 70 + 80 = 260$ cây.
}
\end{ex}

\begin{ex}
Trung bình mỗi khối trồng được bao nhiêu cây?
\shortans{65}
\loigiai{
$260 : 4 = 65$ cây.
}
\end{ex}

\begin{ex}
Cột biểu diễn khối nào cao nhất trong biểu đồ trên? (Khối 4 ghi 4).
\shortans{4}
\loigiai{
Khối 4.
}
\end{ex}

\begin{ex}
Cột biểu diễn khối nào thấp nhất trong biểu đồ trên? (Khối 1 ghi 1).
\shortans{1}
\loigiai{
Khối 1.
}
\end{ex}

\begin{ex}
Biểu đồ cột số học sinh yêu thích các môn thể thao: Bóng đá: 40 bạn, Bóng rổ: 25 bạn, Bơi lội: 35 bạn. Môn thể thao nào được nhiều bạn thích nhất? (Bóng đá ghi 1, Bóng rổ ghi 2, Bơi lội ghi 3).
\shortans{1}
\loigiai{
Bóng đá (40 bạn).
}
\end{ex}

\begin{ex}
Tổng số học sinh tham gia khảo sát là bao nhiêu bạn?
\shortans{100}
\loigiai{
$40 + 25 + 35 = 100$ bạn.
}
\end{ex}

\begin{ex}
Số bạn thích bóng đá nhiều hơn bóng rổ bao nhiêu bạn?
\shortans{15}
\loigiai{
$40 - 25 = 15$ bạn.
}
\end{ex}

\begin{ex}
Biểu đồ số ngày mưa trong các tháng: Tháng 6: 15 ngày, Tháng 7: 18 ngày, Tháng 8: 12 ngày. Tháng nào có nhiều ngày mưa nhất? (Ghi số tháng).
\shortans{7}
\loigiai{
Tháng 7 có 18 ngày mưa.
}
\end{ex}

\begin{ex}
Tổng số ngày mưa trong cả 3 tháng là bao nhiêu ngày?
\shortans{45}
\loigiai{
$15 + 18 + 12 = 45$ ngày.
}
\end{ex}

\begin{ex}
Trung bình mỗi tháng có bao nhiêu ngày mưa?
\shortans{15}
\loigiai{
$45 : 3 = 15$ ngày.
}
\end{ex}

\begin{ex}
Biểu đồ điểm 10 môn Toán của 4 tổ: Tổ 1: 12 điểm, Tổ 2: 15 điểm, Tổ 3: 14 điểm, Tổ 4: 19 điểm. Tổ nào đạt nhiều điểm 10 nhất? (Ghi số tổ).
\shortans{4}
\loigiai{
Tổ 4 đạt 19 điểm 10.
}
\end{ex}

\begin{ex}
Tổ 4 nhiều hơn Tổ 1 bao nhiêu điểm 10?
\shortans{7}
\loigiai{
$19 - 12 = 7$ điểm 10.
}
\end{ex}

\begin{ex}
Tổng số điểm 10 của cả 4 tổ là bao nhiêu?
\shortans{60}
\loigiai{
$12 + 15 + 14 + 19 = 60$ điểm 10.
}
\end{ex}

\begin{ex}
Trung bình mỗi tổ đạt bao nhiêu điểm 10?
\shortans{15}
\loigiai{
$60 : 4 = 15$ điểm 10.
}
\end{ex}

\begin{ex}
Biểu đồ cột gồm có trục nằm ngang và trục gì? (Trục đứng ghi 1, trục xiên ghi 2).
\shortans{1}
\loigiai{
Trục đứng (1).
}
\end{ex}

\begin{ex}
Các cột trong biểu đồ cột có bề rộng thế nào với nhau? (Bằng nhau ghi 1, khác nhau ghi 2).
\shortans{1}
\loigiai{
Bằng nhau (1).
}
\end{ex}

\begin{ex}
Chiều cao của mỗi cột biểu diễn đại lượng nào? (Số liệu ghi 1, tên đối tượng ghi 2).
\shortans{1}
\loigiai{
Số liệu tương ứng (1).
}
\end{ex}

\begin{ex}
Biểu đồ bán xe máy trong 3 tháng: Tháng Một: 120 xe, Tháng Hai: 150 xe, Tháng Ba: 180 xe. Tháng Ba bán nhiều hơn Tháng Một bao nhiêu xe?
\shortans{60}
\loigiai{
$180 - 120 = 60$ xe.
}
\end{ex}

\begin{ex}
Tổng số xe bán trong 3 tháng là bao nhiêu chiếc?
\shortans{450}
\loigiai{
$120 + 150 + 180 = 450$ xe.
}
\end{ex}

\begin{ex}
Trung bình mỗi tháng bán được bao nhiêu chiếc xe?
\shortans{150}
\loigiai{
$450 : 3 = 150$ xe.
}
\end{ex}

\begin{ex}
Nếu cột của tổ 1 cao 12 ô vuông, cột của tổ 2 cao 15 ô vuông. Cột tổ 2 cao hơn tổ 1 bao nhiêu ô vuông?
\shortans{3}
\loigiai{
$15 - 12 = 3$ ô vuông.
}
\end{ex}

\begin{ex}
Khoảng cách giữa các cột trên biểu đồ thường thế nào với nhau? (Cách đều nhau ghi 1, không đều ghi 0).
\shortans{1}
\loigiai{
Cách đều nhau (1).
}
\end{ex}

\begin{ex}
Một biểu đồ cột có 5 cột biểu diễn 5 ngày trong tuần. Có bao nhiêu cột?
\shortans{5}
\loigiai{
Có 5 cột.
}
\end{ex}
